% GENERATED FILE. Edit canonical lessons, then run npm run generate:books.
% canonical-source-hash: fnv1a64:0a51b589b08e05f8
% canonical-lessons: LA-C480-rupes, LA-C480-cautes, LA-C480-praeruptus, LA-C480-arduus, LA-C480-scando

\chapter{Cliff, Jagged rock, Sheer, Steep, To climb}
\label{ch:la-cliff-jagged-rock-sheer-steep-to-climb}
\hlchaptermodality{480}

\begin{chapteropening}
\textbf{By the end of this chapter:} \emph{I can describe mountains: a cliff, a jagged rock, sheer and steep slopes, and climbing.}

Complete the last lesson of chapter 480: I can describe mountains: a cliff, a jagged rock, sheer and steep slopes, and climbing.
\end{chapteropening}

\section[rūpēs, rūpis]{rūpēs, rūpis — a cliff, a rock face}

\label{lesson:LA-C480-rupes}



\begin{quote}
Before the new one: say the Latin for to graft, then the Latin for the grape harvest.
\end{quote}

\subsection*{You'll want to know: rūpēs, rūpis}
\textbf{rūpēs, rūpis} — \textquotedblleft{}a cliff, a rock face\textquotedblright{}.

\begin{quote}
Caprae in rūpibus altīs ambulant.
\end{quote}

\begin{quote}
\emph{Goats walk on the high cliffs.}
\end{quote}

Third declension, feminine: \emph{rūpēs, rūpis f.} A broken rock face, from \emph{rumpō} 'to break'.

\begin{cousinweb}
\emph{rūpēs} is built on \emph{rumpō} 'to break' (Lewis \& Short): rock broken off, a cliff. EtymDB lists it without a parent.
\end{cousinweb}

\begin{grammarlens}[title={What you've built}]
One more word for telling, explaining and describing at B1.
\end{grammarlens}

\subsection*{Your turn}
\begin{itemize}
\raggedright
  \item \emph{Say it:} \emph{rūpēs, rūpis}
  \item \emph{Say it:} \emph{rūpēs, rūpis}, once more
  \item \emph{Say it:} a sentence of your own with \emph{rūpēs, rūpis}
  \item \emph{Recall:} say the Latin for to graft, then the Latin for the grape harvest, then say \emph{rūpēs, rūpis} again
\end{itemize}

\begin{tcolorbox}[breakable,colback=teal!4,colframe=teal!35!black,title={Before you move on}]
What does rūpēs, rūpis mean? (\textquotedblleft{}A cliff, a rock face\textquotedblright{}.) Say it once more.
\end{tcolorbox}
\section[cautēs, cautis]{cautēs, cautis — a jagged rock, a crag}

\label{lesson:LA-C480-cautes}



\begin{quote}
Before the new one: say the Latin for the grape harvest, then the Latin for a cliff.
\end{quote}

\subsection*{You'll want to know: cautēs, cautis}
\textbf{cautēs, cautis} — \textquotedblleft{}a jagged rock, a crag\textquotedblright{}.

\begin{quote}
Nāvis in cautēs cadit et frangitur.
\end{quote}

\begin{quote}
\emph{The ship runs onto the jagged rocks and breaks up.}
\end{quote}

Third declension, feminine: \emph{cautēs, cautis f.}, mostly plural. Sharp rocks on a coast or a mountain.

\begin{cousinweb}
\emph{cautēs} is inherited; Lewis \& Short and EtymDB connect it with \emph{cōs} 'whetstone', from PIE \emph{\'{k}eh\textsubscript{3}-} 'to sharpen': a sharp rock.
\end{cousinweb}

\begin{grammarlens}[title={What you've built}]
One more word for telling, explaining and describing at B1.
\end{grammarlens}

\subsection*{Your turn}
\begin{itemize}
\raggedright
  \item \emph{Say it:} \emph{cautēs, cautis}
  \item \emph{Say it:} \emph{cautēs, cautis}, once more
  \item \emph{Say it:} a sentence of your own with \emph{cautēs, cautis}
  \item \emph{Recall:} say the Latin for the grape harvest, then the Latin for a cliff, then say \emph{cautēs, cautis} again
\end{itemize}

\begin{tcolorbox}[breakable,colback=teal!4,colframe=teal!35!black,title={Before you move on}]
What does cautēs, cautis mean? (\textquotedblleft{}A jagged rock, a crag\textquotedblright{}.) Say it once more.
\end{tcolorbox}
\section[praeruptus, praerupta, praeruptum]{praeruptus, praerupta, praeruptum — sheer, very steep, broken off}

\label{lesson:LA-C480-praeruptus}



\begin{quote}
Before the new one: say the Latin for a cliff, then the Latin for a jagged rock.
\end{quote}

\subsection*{You'll want to know: praeruptus, praerupta, praeruptum}
\textbf{praeruptus, praerupta, praeruptum} — \textquotedblleft{}sheer, very steep, broken off\textquotedblright{}.

\begin{quote}
Via in monte praerupta est, itaque mūlī lentē eunt.
\end{quote}

\begin{quote}
\emph{The road up the mountain is very steep, so the mules go slowly.}
\end{quote}

An adjective of three endings, the participle of \emph{praerumpō} 'to break off in front'. Of rock: sheer.

\begin{cousinweb}
\emph{praeruptus} is the past participle of \emph{praerumpō} 'to break off in front', from \emph{prae-} and \emph{rumpō} 'to break' (Lewis \& Short): land broken off sharply.
\end{cousinweb}

\begin{grammarlens}[title={What you've built}]
One more word for telling, explaining and describing at B1.
\end{grammarlens}

\subsection*{Your turn}
\begin{itemize}
\raggedright
  \item \emph{Say it:} \emph{praeruptus, praerupta, praeruptum}
  \item \emph{Say it:} \emph{praeruptus, praerupta, praeruptum}, once more
  \item \emph{Say it:} a sentence of your own with \emph{praeruptus, praerupta, praeruptum}
  \item \emph{Recall:} say the Latin for a cliff, then the Latin for a jagged rock, then say \emph{praeruptus, praerupta, praeruptum} again
\end{itemize}

\begin{tcolorbox}[breakable,colback=teal!4,colframe=teal!35!black,title={Before you move on}]
What does praeruptus, praerupta, praeruptum mean? (\textquotedblleft{}Sheer, very steep, broken off\textquotedblright{}.) Say it once more.
\end{tcolorbox}
\section[arduus, ardua, arduum]{arduus, ardua, arduum — steep, high; difficult}

\label{lesson:LA-C480-arduus}



\begin{quote}
Before the new one: say the Latin for a jagged rock, then the Latin for sheer.
\end{quote}

\subsection*{You'll want to know: arduus, ardua, arduum}
\textbf{arduus, ardua, arduum} — \textquotedblleft{}steep, high; difficult\textquotedblright{}.

\begin{quote}
Sēmita ardua est, et viātōrēs lassī sunt.
\end{quote}

\begin{quote}
\emph{The path is steep, and the travellers are weary.}
\end{quote}

An adjective of three endings. \emph{Rēs ardua}, a hard task. English \emph{arduous} comes from it.

\begin{cousinweb}
\emph{arduus} is inherited, but its origin is not settled: Lewis \& Short link it, doubtfully, with \emph{alō} 'to nourish' and \emph{altus} 'high', and EtymDB gives no parent. English \emph{arduous} comes from it.
\end{cousinweb}

\begin{grammarlens}[title={What you've built}]
One more word for telling, explaining and describing at B1.
\end{grammarlens}

\subsection*{Your turn}
\begin{itemize}
\raggedright
  \item \emph{Say it:} \emph{arduus, ardua, arduum}
  \item \emph{Say it:} \emph{arduus, ardua, arduum}, once more
  \item \emph{Say it:} a sentence of your own with \emph{arduus, ardua, arduum}
  \item \emph{Recall:} say the Latin for a jagged rock, then the Latin for sheer, then say \emph{arduus, ardua, arduum} again
\end{itemize}

\begin{tcolorbox}[breakable,colback=teal!4,colframe=teal!35!black,title={Before you move on}]
What does arduus, ardua, arduum mean? (\textquotedblleft{}Steep, high; difficult\textquotedblright{}.) Say it once more.
\end{tcolorbox}
\section[scandō, scandere]{scandō, scandere — to climb, to mount}

\label{lesson:LA-C480-scando}



\begin{quote}
Before the new one: say the Latin for sheer, then the Latin for steep.
\end{quote}

\subsection*{You'll want to know: scandō, scandere}
\textbf{scandō, scandere} — \textquotedblleft{}to climb, to mount\textquotedblright{}.

\begin{quote}
Amīcī montem scandunt et urbem procul cōnspicantur.
\end{quote}

\begin{quote}
\emph{The friends climb the mountain and catch sight of the city far off.}
\end{quote}

\emph{scandō, scandere, scandī, scānsum}. English \emph{scan} (to count out the feet of a verse) comes from it, and \emph{ascend} and \emph{descend} from its compounds.

\begin{cousinweb}
\emph{scandō} is inherited from PIE \emph{skend-} 'to jump, to climb' (EtymDB). English \emph{scan} comes from it, and \emph{ascend} and \emph{descend} from its compounds \emph{ascendō} and \emph{dēscendō} (Lewis \& Short).
\end{cousinweb}

\begin{grammarlens}[title={What you've built}]
One more word for telling, explaining and describing at B1.
\end{grammarlens}

\subsection*{Your turn}
\begin{itemize}
\raggedright
  \item \emph{Say it:} \emph{scandō, scandere}
  \item \emph{Say it:} \emph{scandō, scandere}, once more
  \item \emph{Say it:} a sentence of your own with \emph{scandō, scandere}
  \item \emph{Recall:} say the Latin for sheer, then the Latin for steep, then say \emph{scandō, scandere} again
\end{itemize}

\begin{tcolorbox}[breakable,colback=teal!4,colframe=teal!35!black,title={Before you move on}]
What does scandō, scandere mean? (\textquotedblleft{}To climb, to mount\textquotedblright{}.) Say it once more.
\end{tcolorbox}
